This study reexamines the Standard Least-Squares Criterion (SLSC) used in hydrological frequency analysis, focusing on sample-size dependence and the effect of standardization. Monte Carlo simulations for six distributions show that SLSC tends to decrease as sample size increases. By distinguishing the original data space (X-space) from the standardized plotting-paper space (S-space), the analysis shows that nonlinear standardization makes SLSC a distance measure that depends on the standardizing function. In particular, the S-space evaluation of LN3 tends to give small SLSC values and can affect comparisons among distributions. The outcomes of the SLSC+jackknife procedure and AIC were also compared. These results indicate that SLSC should be interpreted with attention to fixed thresholds, the definition of the standard variate, and the resulting distance scale.
